\documentclass[11pt,a4paper]{article}
\usepackage{shared/parscover}

% ============================================================================
% Post-Quantum Cryptographic Standards for Pars Network
% Pars Network Technical Whitepaper PNP-004
% ============================================================================

\usepackage[utf8]{inputenc}
\usepackage[T1]{fontenc}
\usepackage{amsmath,amssymb,amsthm}
\usepackage{graphicx}
\usepackage{hyperref}
\usepackage{algorithm}
\usepackage{algpseudocode}
\usepackage{booktabs}
\usepackage{listings}
\input{shared/lstlang}
\usepackage{xcolor}
\usepackage{tikz}
\usepackage{enumitem}
\usepackage{caption}
\usepackage{fancyhdr}
\usepackage{abstract}

\geometry{margin=1in}

\hypersetup{
    colorlinks=true,
    linkcolor=blue!70!black,
    citecolor=green!50!black,
    urlcolor=blue!60!black
}

\lstset{
    basicstyle=\ttfamily\small,
    breaklines=true,
    frame=single,
    backgroundcolor=\color{gray!10},
    keywordstyle=\color{blue!70!black}\bfseries,
    commentstyle=\color{green!50!black}\itshape,
    stringstyle=\color{red!60!black},
    numbers=left,
    numberstyle=\tiny\color{gray},
    numbersep=5pt
}

\usetikzlibrary{arrows.meta,positioning,shapes.geometric,fit,calc}

\newtheorem{definition}{Definition}[section]
\newtheorem{theorem}{Theorem}[section]
\newtheorem{lemma}{Lemma}[section]
\newtheorem{proposition}{Proposition}[section]
\newtheorem{corollary}{Corollary}[section]

\pagestyle{fancy}
\fancyhf{}
\fancyhead[L]{\small Cyrus Pahlavi, Pars Team}
\fancyhead[R]{\small Post-Quantum Cryptography}
\fancyfoot[C]{\thepage}

% ============================================================================
\title{%
    \textbf{Post-Quantum Cryptographic Standards\\
    for Pars Network}\\[0.5em]
    \large Pars Network Technical Whitepaper PNP-004
}

\author{%
    Cyrus Pahlavi, Pars Team\\
    \texttt{research@pars.network}\\[0.5em]
    \small Version 1.0 --- February 2026
}

\date{}

% ============================================================================
\begin{document}
\parscoverpage

\thispagestyle{fancy}

% ============================================================================
\begin{abstract}
\noindent
We present the post-quantum cryptographic migration strategy for the Pars
Network, addressing the ``harvest now, decrypt later'' threat relevant to
diaspora communities under state surveillance.  Adversaries may record
encrypted traffic today for decryption by future quantum computers.  We
specify hybrid schemes combining classical primitives (X25519, Ed25519)
with NIST-standardised post-quantum algorithms: ML-KEM (FIPS~203) for key
encapsulation, ML-DSA (FIPS~204) for digital signatures, and SLH-DSA
(FIPS~205) for stateless hash-based signatures.  A phased migration plan
maintains backward compatibility while strengthening all layers.  Benchmarks
on mobile hardware demonstrate that hybrid schemes add less than 3\,ms to
session establishment and less than 2\,KB to signature sizes.  We provide
parameter selections, implementation guidance, and formal analysis of hybrid
composition security under the assumption that at least one component
remains secure.
\end{abstract}

\vspace{1em}
\noindent\textbf{Keywords:} post-quantum cryptography, ML-KEM, ML-DSA,
SLH-DSA, hybrid schemes, lattice-based, hash-based, migration strategy

\tableofcontents
\newpage

% ============================================================================
\section{Introduction}
\label{sec:introduction}

Cryptographically relevant quantum computers (CRQCs) threaten the
foundations of all blockchain and communication systems.  Shor's
algorithm~\cite{shor1994algorithms} breaks RSA, DSA, and ECDH/ECDSA in
polynomial time.  Grover's algorithm~\cite{grover1996fast} halves the
effective security of symmetric ciphers and hash functions.

For diaspora communities under state surveillance, the ``harvest now,
decrypt later'' (HNDL) strategy---where adversaries record encrypted
communications for future quantum decryption---is a documented
practice~\cite{nsareport2016quantum}.  Traffic encrypted today with
classical-only schemes may become readable within 10--20~years.

\subsection{Threat Timeline}

\begin{table}[h]
\centering
\caption{Quantum Threat Timeline Estimates}
\label{tab:timeline}
\begin{tabular}{@{}lll@{}}
\toprule
\textbf{Estimate} & \textbf{Source} & \textbf{CRQC Year} \\
\midrule
Optimistic   & Mosca~\cite{mosca2018cybersecurity} & 2030--2035 \\
Median       & NIST~\cite{nist2016pqc}             & 2035--2040 \\
Conservative & NSA CNSA~2.0~\cite{nsa2022cnsa2}    & 2040--2050 \\
\bottomrule
\end{tabular}
\end{table}

Given 10--15~year confidentiality requirements and 5--10~year migration
timelines, the window for action is now.

\subsection{Contributions}

\begin{itemize}
    \item Concrete parameter selections for all Pars post-quantum
    primitives.
    \item Hybrid composition framework with formal security reduction.
    \item Phased migration preserving backward compatibility.
    \item Performance benchmarks on mobile and embedded hardware.
    \item Integration specifications for the Pars L2 (PNP-001), session
    protocol (PNP-003), and DID system (PNP-008).
\end{itemize}

% ============================================================================
\section{Background}
\label{sec:background}

\subsection{NIST Post-Quantum Standards}

NIST finalised three standards in 2024~\cite{nist2024pqc}:

\begin{enumerate}
    \item \textbf{ML-KEM (FIPS~203):} Module-Lattice Key Encapsulation,
    based on CRYSTALS-Kyber~\cite{bos2018crystals}.
    \item \textbf{ML-DSA (FIPS~204):} Module-Lattice Digital Signature,
    based on CRYSTALS-Dilithium~\cite{ducas2018crystals}.
    \item \textbf{SLH-DSA (FIPS~205):} Stateless Hash-Based Signatures,
    based on SPHINCS+~\cite{bernstein2019sphincs}.
\end{enumerate}

\subsection{Hardness Assumptions}

\begin{definition}[Module-LWE]
    Given $\mathbf{A} \in \mathbb{Z}_q^{k \times k}$ and secret
    $\mathbf{s} \in \mathbb{Z}_q^k$ with small entries, distinguish
    $(\mathbf{A}, \mathbf{As} + \mathbf{e})$ from
    $(\mathbf{A}, \mathbf{u})$ where $\mathbf{e}$ is small error and
    $\mathbf{u}$ is uniform.
\end{definition}

\begin{definition}[Module-SIS]
    Given $\mathbf{A} \in \mathbb{Z}_q^{k \times \ell}$, find non-zero
    $\mathbf{z} \in \mathbb{Z}^{\ell}$ with
    $\|\mathbf{z}\|_\infty \leq \beta$ and
    $\mathbf{Az} = \mathbf{0} \bmod q$.
\end{definition}

ML-KEM relies on M-LWE; ML-DSA on both M-LWE and M-SIS.  SLH-DSA relies
only on hash function security (second-preimage resistance).

\subsection{Current Cryptographic Inventory}

\begin{table}[h]
\centering
\caption{Current Primitives and Quantum Vulnerability}
\label{tab:inventory}
\begin{tabular}{@{}llll@{}}
\toprule
\textbf{Usage} & \textbf{Primitive} & \textbf{Classical} & \textbf{PQ Status} \\
\midrule
Key agreement  & X25519      & 128-bit & Broken (Shor) \\
Signatures     & Ed25519     & 128-bit & Broken (Shor) \\
Encryption     & AES-256-GCM & 256-bit & 128-bit (Grover) \\
Hashing        & SHA-256     & 256-bit & 128-bit (Grover) \\
Sessions       & X3DH + DR   & 128-bit & Broken (DH) \\
L2 accounts    & secp256k1   & 128-bit & Broken (Shor) \\
\bottomrule
\end{tabular}
\end{table}

% ============================================================================
\section{Algorithm Selection}
\label{sec:selection}

\subsection{ML-KEM for Key Encapsulation}

\begin{table}[h]
\centering
\caption{ML-KEM Parameter Sets}
\label{tab:mlkem}
\begin{tabular}{@{}lrrrrr@{}}
\toprule
\textbf{Level} & $k$ & \textbf{PK} & \textbf{CT} & \textbf{SS} &
    \textbf{NIST Lvl} \\
\midrule
ML-KEM-512  & 2 &  800\,B &  768\,B & 32\,B & 1 \\
ML-KEM-768  & 3 & 1184\,B & 1088\,B & 32\,B & 3 \\
ML-KEM-1024 & 4 & 1568\,B & 1568\,B & 32\,B & 5 \\
\bottomrule
\end{tabular}
\end{table}

\textbf{Selection: ML-KEM-768.}  Level~3 matches 192-bit classical security.
Key/ciphertext sizes below 1.2\,KB.  Operations under 0.5\,ms on mobile ARM.

\subsection{ML-DSA for Digital Signatures}

\begin{table}[h]
\centering
\caption{ML-DSA Parameter Sets}
\label{tab:mldsa}
\begin{tabular}{@{}lrrrr@{}}
\toprule
\textbf{Level} & \textbf{PK} & \textbf{SK} & \textbf{Sig} &
    \textbf{NIST Lvl} \\
\midrule
ML-DSA-44 & 1312\,B & 2560\,B & 2420\,B & 2 \\
ML-DSA-65 & 1952\,B & 4032\,B & 3309\,B & 3 \\
ML-DSA-87 & 2592\,B & 4896\,B & 4627\,B & 5 \\
\bottomrule
\end{tabular}
\end{table}

\textbf{Selection: ML-DSA-65.}  Level~3 matches ML-KEM-768.  Signature
(3.3\,KB) acceptable for L2 transactions.  Signing $<1$\,ms.

\subsection{SLH-DSA for Root-of-Trust}

\begin{table}[h]
\centering
\caption{SLH-DSA Parameter Sets (Selected)}
\label{tab:slhdsa}
\begin{tabular}{@{}lrrrr@{}}
\toprule
\textbf{Variant} & \textbf{PK} & \textbf{SK} & \textbf{Sig} &
    \textbf{NIST Lvl} \\
\midrule
SLH-DSA-SHA2-128s & 32\,B  & 64\,B  &  7856\,B & 1 \\
SLH-DSA-SHA2-192f & 48\,B  & 96\,B  & 35664\,B & 3 \\
SLH-DSA-SHA2-256s & 64\,B  & 128\,B & 29792\,B & 5 \\
\bottomrule
\end{tabular}
\end{table}

\textbf{Selection: SLH-DSA-SHA2-128s} for genesis block, protocol
parameters, and governance keys.  Security rests solely on hash functions
(insurance against lattice breakthroughs).  Used infrequently, so larger
signatures are acceptable.

% ============================================================================
\section{Hybrid Composition Framework}
\label{sec:hybrid}

\subsection{Design Principle}

\begin{definition}[Hybrid Security]
    Scheme $\mathcal{H} = (\mathcal{C}, \mathcal{P})$ with classical
    $\mathcal{C}$ and post-quantum $\mathcal{P}$ satisfies:
    \[
        \mathrm{Adv}_{\mathcal{H}}^{\mathrm{IND}} \leq
        \min\bigl(\mathrm{Adv}_{\mathcal{C}}^{\mathrm{IND}},\;
        \mathrm{Adv}_{\mathcal{P}}^{\mathrm{IND}}\bigr)
        + \mathrm{negl}(\lambda).
    \]
\end{definition}

\subsection{Hybrid Key Encapsulation}

\begin{algorithm}[H]
\caption{Hybrid KEM: Encapsulate}
\label{alg:hybrid-enc}
\begin{algorithmic}[1]
\Require Classical PK $\mathrm{pk}_C$ (X25519), PQ PK $\mathrm{pk}_P$
    (ML-KEM-768)
\State $(x, X) \gets$ fresh X25519 pair
\State $\mathrm{ss}_C \gets \mathrm{X25519}(x, \mathrm{pk}_C)$
\State $(\mathrm{ct}_P, \mathrm{ss}_P) \gets
    \mathrm{ML\text{-}KEM.Enc}(\mathrm{pk}_P)$
\State $\mathrm{SS} \gets \mathrm{HKDF}\bigl(
    \mathrm{ss}_C \| \mathrm{ss}_P,\;
    X \| \mathrm{ct}_P,\;
    \texttt{"Pars-HybridKEM-v1"},\; 32\bigr)$
\State \Return $(X \| \mathrm{ct}_P,\;\mathrm{SS})$
\end{algorithmic}
\end{algorithm}

\begin{algorithm}[H]
\caption{Hybrid KEM: Decapsulate}
\label{alg:hybrid-dec}
\begin{algorithmic}[1]
\Require $\mathrm{sk}_C$, $\mathrm{sk}_P$, ciphertext $(X, \mathrm{ct}_P)$
\State $\mathrm{ss}_C \gets \mathrm{X25519}(\mathrm{sk}_C, X)$
\State $\mathrm{ss}_P \gets
    \mathrm{ML\text{-}KEM.Dec}(\mathrm{sk}_P, \mathrm{ct}_P)$
\State $\mathrm{SS} \gets \mathrm{HKDF}\bigl(
    \mathrm{ss}_C \| \mathrm{ss}_P,\;
    X \| \mathrm{ct}_P,\;
    \texttt{"Pars-HybridKEM-v1"},\; 32\bigr)$
\State \Return SS
\end{algorithmic}
\end{algorithm}

\begin{theorem}[Hybrid KEM IND-CCA2]
    \label{thm:hybrid-kem}
    The hybrid KEM is IND-CCA2 secure if either X25519 (Gap-CDH) or
    ML-KEM-768 (M-LWE) is IND-CCA2 secure.
\end{theorem}

\begin{proof}
    The shared secret $\mathrm{SS} = \mathrm{HKDF}(\mathrm{ss}_C \|
    \mathrm{ss}_P, \ldots)$.  By the HKDF extraction
    lemma~\cite{krawczyk2010hkdf}, if either input is uniformly random and
    independent of the adversary's view, the output is computationally
    random.  Under Gap-CDH, $\mathrm{ss}_C$ is random.  Under M-LWE,
    $\mathrm{ss}_P$ is random.  Since HKDF requires only one random
    component, the adversary's advantage is negligible unless both
    components are broken.
\end{proof}

\subsection{Hybrid Digital Signatures}

\begin{definition}[Hybrid Signature]
    On message $m$:
    \[
        \sigma_H = \bigl(
            \mathrm{Ed25519.Sign}(\mathrm{sk}_C, m),\;
            \mathrm{ML\text{-}DSA.Sign}(\mathrm{sk}_P, m)
        \bigr).
    \]
    Verification requires both components to pass.
\end{definition}

\begin{theorem}[Hybrid Signature EUF-CMA]
    The hybrid signature is EUF-CMA secure if either Ed25519 or ML-DSA-65
    is EUF-CMA secure.
\end{theorem}

\begin{proof}
    A forger must produce both valid $\sigma_C$ and $\sigma_P$.  If
    Ed25519 is secure, $\sigma_C$ cannot be forged.  If ML-DSA is
    secure, $\sigma_P$ cannot be forged.  The forger succeeds only if
    both are broken simultaneously.
\end{proof}

% ============================================================================
\section{Migration Strategy}
\label{sec:migration}

\subsection{Phased Approach}

\begin{table}[h]
\centering
\caption{Migration Phases}
\label{tab:phases}
\begin{tabular}{@{}llll@{}}
\toprule
\textbf{Phase} & \textbf{Timeline} & \textbf{Scope} & \textbf{Mode} \\
\midrule
0: Preparation & Q1 2026 & Library integration  & Classical only \\
1: Hybrid opt  & Q2 2026 & Session protocol     & Hybrid (optional) \\
2: Hybrid dflt & Q3 2026 & All new sessions     & Hybrid (default) \\
3: Mandatory   & Q1 2027 & All connections      & Hybrid (required) \\
4: PQ primary  & 2028+   & Deprecate classical  & PQ (with fallback) \\
\bottomrule
\end{tabular}
\end{table}

\subsection{Phase~0: Library Integration}

\begin{lstlisting}[language=Go,caption={PQ Crypto Interface}]
type PQCrypto interface {
    KEMKeyGen() (KEMPub, KEMSec, error)
    KEMEncapsulate(pk KEMPub) (
        ct, ss []byte, err error,
    )
    KEMDecapsulate(
        sk KEMSec, ct []byte,
    ) ([]byte, error)

    DSAKeyGen() (DSAPub, DSASec, error)
    DSASign(sk DSASec, msg []byte) (
        []byte, error,
    )
    DSAVerify(
        pk DSAPub, msg, sig []byte,
    ) (bool, error)
}

type MLKEM768 struct{}

func (m *MLKEM768) KEMKeyGen() (
    KEMPub, KEMSec, error,
) {
    pk, sk, err := mlkem768.GenerateKey()
    if err != nil {
        return nil, nil,
            fmt.Errorf("keygen: %w", err)
    }
    return pk, sk, nil
}

func (m *MLKEM768) KEMEncapsulate(
    pk KEMPub,
) ([]byte, []byte, error) {
    ct, ss, err := mlkem768.Encapsulate(pk)
    if err != nil {
        return nil, nil,
            fmt.Errorf("encap: %w", err)
    }
    return ct, ss, nil
}

func (m *MLKEM768) KEMDecapsulate(
    sk KEMSec, ct []byte,
) ([]byte, error) {
    ss, err := mlkem768.Decapsulate(sk, ct)
    if err != nil {
        return nil,
            fmt.Errorf("decap: %w", err)
    }
    return ss, nil
}
\end{lstlisting}

\subsection{Phase~1: Hybrid Session Protocol}

PSP session establishment (PNP-003) extended with hybrid X3DH.  Prekey
bundles include PQ public keys.  Nodes advertise PQ capability via
version flag.

\subsection{Phase~2--3: Default and Mandatory}

All new sessions default to hybrid (Phase~2).  Classical-only rejected
(Phase~3).  All L2 transactions require hybrid signatures.

\subsection{Phase~4: PQ Primary}

Classical keys maintained only for external interoperability.  Primary
security rests on PQ algorithms.

\subsection{L2 Account Migration}

\begin{lstlisting}[language=Solidity,caption={Hybrid Account Registry}]
contract HybridAccountRegistry {
    struct HybridAccount {
        address classicalAddr;
        bytes pqPubKey;       // ML-DSA-65
        bytes32 rootKeyHash;  // SLH-DSA
        uint8 phase;
        uint256 pqRegistered;
    }

    mapping(address => HybridAccount) public accts;

    function registerPQKey(
        bytes calldata pqPub,
        bytes calldata pqSig,
        bytes calldata classSig
    ) external {
        require(
            _verifyClassical(
                msg.sender, pqPub, classSig
            ), "Bad classical sig"
        );
        require(
            _verifyMLDSA(
                pqPub,
                abi.encodePacked(
                    msg.sender, pqPub
                ),
                pqSig
            ), "Bad PQ sig"
        );
        accts[msg.sender] = HybridAccount({
            classicalAddr: msg.sender,
            pqPubKey: pqPub,
            rootKeyHash: bytes32(0),
            phase: 1,
            pqRegistered: block.timestamp
        });
    }
}
\end{lstlisting}

% ============================================================================
\section{Performance Analysis}
\label{sec:performance}

\subsection{Operation Timing}

\begin{table}[h]
\centering
\caption{PQ Operation Timing (mobile ARM)}
\label{tab:pq-perf}
\begin{tabular}{@{}lrr@{}}
\toprule
\textbf{Operation} & \textbf{A17} & \textbf{SD 8g3} \\
\midrule
ML-KEM-768 KeyGen       & 0.08\,ms & 0.12\,ms \\
ML-KEM-768 Encapsulate  & 0.11\,ms & 0.16\,ms \\
ML-KEM-768 Decapsulate  & 0.09\,ms & 0.14\,ms \\
ML-DSA-65 KeyGen        & 0.15\,ms & 0.22\,ms \\
ML-DSA-65 Sign          & 0.52\,ms & 0.71\,ms \\
ML-DSA-65 Verify        & 0.18\,ms & 0.26\,ms \\
SLH-DSA-SHA2-128s Sign  & 142\,ms  & 198\,ms  \\
SLH-DSA-SHA2-128s Verify& 6.2\,ms  & 8.8\,ms  \\
\midrule
Hybrid KEM              & 0.19\,ms & 0.28\,ms \\
Hybrid Sig              & 0.57\,ms & 0.76\,ms \\
\bottomrule
\end{tabular}
\end{table}

\subsection{Size Overhead}

\begin{table}[h]
\centering
\caption{Key and Signature Sizes}
\label{tab:sizes}
\begin{tabular}{@{}lrrr@{}}
\toprule
\textbf{Scheme} & \textbf{PK} & \textbf{SK} & \textbf{Sig/CT} \\
\midrule
X25519 (classical)  & 32\,B   & 32\,B   & 32\,B \\
ML-KEM-768          & 1184\,B & 2400\,B & 1088\,B \\
Hybrid KEM          & 1216\,B & 2432\,B & 1120\,B \\
\midrule
Ed25519 (classical) & 32\,B   & 64\,B   & 64\,B \\
ML-DSA-65           & 1952\,B & 4032\,B & 3309\,B \\
Hybrid Sig          & 1984\,B & 4096\,B & 3373\,B \\
\midrule
SLH-DSA-SHA2-128s   & 32\,B   & 64\,B   & 7856\,B \\
\bottomrule
\end{tabular}
\end{table}

\subsection{Session Establishment Impact}

Hybrid X3DH adds 1.3\,ms over classical-only (2.1\,ms vs 0.8\,ms),
negligible relative to network latency (50--500\,ms).

\subsection{L2 Transaction Size Impact}

\begin{table}[h]
\centering
\caption{Transaction Size Impact}
\label{tab:tx-size}
\begin{tabular}{@{}lrr@{}}
\toprule
\textbf{Component} & \textbf{Classical} & \textbf{Hybrid} \\
\midrule
Addresses + value    & 72\,B    & 72\,B \\
Classical signature  & 65\,B    & 65\,B \\
PQ signature         & ---      & 3309\,B \\
\midrule
\textbf{Total}       & \textbf{137\,B} & \textbf{3446\,B} \\
\bottomrule
\end{tabular}
\end{table}

Mitigations: (1)~ML-DSA batch aggregation via Fiat-Shamir;
(2)~$\approx 40\%$ zlib compression; (3)~in rollup model, PQ sigs are
verified by L2 validators, not posted to L1.

% ============================================================================
\section{Implementation Considerations}
\label{sec:implementation}

\subsection{Side-Channel Resistance}

All PQ implementations must be constant-time:

\begin{lstlisting}[language=Go,caption={Constant-Time Utilities}]
func ctEqual(a, b []byte) bool {
    if len(a) != len(b) {
        return false
    }
    var acc byte
    for i := range a {
        acc |= a[i] ^ b[i]
    }
    return acc == 0
}

func ctSelect(
    cond bool, a, b []byte,
) []byte {
    mask := byte(0)
    if cond {
        mask = 0xFF
    }
    out := make([]byte, len(a))
    for i := range a {
        out[i] = (a[i] & mask) |
                 (b[i] & ^mask)
    }
    return out
}
\end{lstlisting}

ML-KEM decapsulation uses implicit rejection (FO transform): on
decryption failure, return a pseudorandom SS derived from secret key and
ciphertext rather than an error, preventing chosen-ciphertext attacks.

\subsection{Random Number Generation}

Requirements:
\begin{enumerate}
    \item OS CSPRNG (\texttt{/dev/urandom}, \texttt{SecRandomCopyBytes}).
    \item Hardware RNG mixing where available.
    \item NIST SP~800-90A DRBG seeded from OS entropy.
\end{enumerate}

\subsection{Key Storage}

\begin{table}[h]
\centering
\caption{Key Storage Requirements}
\label{tab:storage}
\begin{tabular}{@{}lrr@{}}
\toprule
\textbf{Key Type} & \textbf{Size} & \textbf{Location} \\
\midrule
ML-KEM-768 pair     & 3.6\,KB  & Secure enclave \\
ML-DSA-65 pair      & 6.0\,KB  & Secure enclave \\
SLH-DSA root pair   & 0.2\,KB  & Hardware token \\
Prekey bundle       & 2.5\,KB  & Session directory \\
100 hybrid OTPKs    & 125\,KB  & Session directory \\
\bottomrule
\end{tabular}
\end{table}

% ============================================================================
\section{Security Analysis}
\label{sec:security}

\subsection{Post-Migration Security Levels}

\begin{table}[h]
\centering
\caption{Security After Migration}
\label{tab:security-post}
\begin{tabular}{@{}llll@{}}
\toprule
\textbf{Usage} & \textbf{Primitive} & \textbf{Classical} & \textbf{Quantum} \\
\midrule
Key agreement & Hybrid KEM   & 128-bit & $\geq$128-bit \\
Tx signatures & Hybrid Sig   & 128-bit & $\geq$128-bit \\
Root-of-trust & SLH-DSA-128s & 128-bit & 128-bit \\
Encryption    & AES-256-GCM  & 256-bit & 128-bit \\
Hashing       & SHA-256      & 256-bit & 128-bit \\
\bottomrule
\end{tabular}
\end{table}

\subsection{HNDL Mitigation}

\begin{theorem}[HNDL Resistance]
    Traffic encrypted after Phase~1 resists HNDL if at the time of
    quantum decryption either (a)~X25519 remains secure or (b)~ML-KEM-768
    remains secure.
\end{theorem}

\begin{proof}
    Direct from Theorem~\ref{thm:hybrid-kem}.
\end{proof}

\subsection{Lattice Attack Landscape}

\begin{enumerate}
    \item Current best lattice reduction: BKZ-$\beta$ with $\beta \approx
    400$.  Breaking ML-KEM-768 requires $\beta >
    800$~\cite{albrecht2019estimate}.
    \item Known quantum speedups for lattice sieving are polynomial, not
    exponential~\cite{laarhoven2015quantum}.
    \item Implementation-level side-channel attacks remain the primary
    practical concern.
\end{enumerate}

\subsection{Diversity Argument for SLH-DSA}

SLH-DSA's security rests on hash function properties only, not lattice
assumptions.  If a polynomial-time lattice algorithm is discovered
(breaking ML-KEM and ML-DSA), SLH-DSA-protected root keys remain secure.
This provides cryptographic diversity across the trust hierarchy.

% ============================================================================
\section{Interoperability}
\label{sec:interop}

\subsection{TLS~1.3 Integration}

\begin{lstlisting}[language=Go,caption={Hybrid TLS Configuration}]
cfg := &tls.Config{
    CurvePreferences: []tls.CurveID{
        tls.X25519MLKEM768,
        tls.X25519,
    },
    MinVersion: tls.VersionTLS13,
}
\end{lstlisting}

\subsection{Ethereum Address Compatibility}

Preserve 20-byte address format:
\[
    \mathrm{addr} = \mathrm{keccak256}\bigl(
        \mathrm{pk}_C \| H(\mathrm{pk}_P)
    \bigr)[12:]
\]

\subsection{Cross-Chain Bridge Signatures}

Bridge contracts on Ethereum mainnet require hybrid verification:

\begin{lstlisting}[language=Solidity,caption={Hybrid Verification Precompile}]
// Proposed precompile at 0x0110
// Input: pk_classical || pk_pq ||
//        sig_classical || sig_pq || msg
// Output: bool
function verifyHybrid(
    bytes calldata input
) external view returns (bool) {
    // Decode components
    (bytes memory pkC, bytes memory pkP,
     bytes memory sigC, bytes memory sigP,
     bytes memory msg) = abi.decode(
        input,
        (bytes,bytes,bytes,bytes,bytes)
    );
    // Both must verify
    return ecrecover_verify(pkC, msg, sigC)
        && mldsa_verify(pkP, msg, sigP);
}
\end{lstlisting}

% ============================================================================
\section{Formal Verification}
\label{sec:formal}

\subsection{Implementation Verification}

We target formal verification of:

\begin{enumerate}
    \item \textbf{Constant-time execution:} Using ct-verif and
    Jasmin~\cite{almeida2017jasmin} to verify absence of secret-dependent
    branches and memory accesses.
    \item \textbf{Functional correctness:} Proving decapsulation inverts
    encapsulation and signature verification accepts valid signatures,
    using the F* proof assistant.
    \item \textbf{Protocol composition:} Verifying the hybrid composition
    preserves security using CryptoVerif~\cite{blanchet2008cryptoverif}.
\end{enumerate}

\subsection{Conformance Testing}

All implementations must pass:

\begin{enumerate}
    \item NIST Known-Answer Test (KAT) vectors for ML-KEM-768, ML-DSA-65,
    and SLH-DSA-SHA2-128s.
    \item Pars-specific hybrid KEM and hybrid signature test vectors.
    \item Interoperability tests against reference implementations.
\end{enumerate}

% ============================================================================
\section{Conclusion}
\label{sec:conclusion}

The Pars Network post-quantum migration addresses the immediate HNDL
threat through hybrid schemes composing classical and PQ primitives.
The phased migration preserves compatibility while strengthening all
layers.  Hybrid operations add less than 3\,ms to session establishment.
SLH-DSA provides insurance against lattice breakthroughs for root-of-trust
operations.  Mandatory hybrid mode is targeted for Q1~2027, with
PQ-primary for 2028 and beyond.

% ============================================================================
\bibliographystyle{plain}
\begin{thebibliography}{26}

\bibitem{shor1994algorithms}
P.~Shor,
``Algorithms for quantum computation,''
\emph{FOCS}, pp.~124--134, 1994.

\bibitem{grover1996fast}
L.~Grover,
``A fast quantum mechanical algorithm for database search,''
\emph{STOC}, pp.~212--219, 1996.

\bibitem{nsareport2016quantum}
NSA,
``Commercial National Security Algorithm Suite 2.0,''
2022.

\bibitem{mosca2018cybersecurity}
M.~Mosca,
``Cybersecurity in an Era with Quantum Computers,''
\emph{IEEE Security \& Privacy}, 16(5), 2018.

\bibitem{nist2016pqc}
NIST,
``Post-Quantum Cryptography: Call for Proposals,''
2016.

\bibitem{nsa2022cnsa2}
NSA,
``CNSA 2.0,''
2022.

\bibitem{nist2024pqc}
NIST,
``FIPS 203, 204, 205,''
2024.

\bibitem{bos2018crystals}
J.~Bos \emph{et al.},
``CRYSTALS -- Kyber,''
\emph{Euro S\&P}, 2018.

\bibitem{ducas2018crystals}
L.~Ducas \emph{et al.},
``CRYSTALS -- Dilithium,''
\emph{TCHES}, 2018(1), 2018.

\bibitem{bernstein2019sphincs}
D.\,J.~Bernstein \emph{et al.},
``SPHINCS+,''
\emph{CHES}, 2019.

\bibitem{krawczyk2010hkdf}
H.~Krawczyk,
``HKDF,''
\emph{CRYPTO 2010}, pp.~631--648, 2010.

\bibitem{stebila2022hybrid}
D.~Stebila, S.~Fluhrer, and S.~Gueron,
``Hybrid Key Exchange in TLS 1.3,''
\emph{IETF Draft}, 2022.

\bibitem{albrecht2019estimate}
M.~Albrecht \emph{et al.},
``Estimate All the LWE, NTRU Schemes!,''
\emph{SCN 2018}, 2018.

\bibitem{laarhoven2015quantum}
T.~Laarhoven, M.~Mosca, and J.~van~de~Pol,
``Finding shortest lattice vectors faster using quantum search,''
\emph{Des.\ Codes Cryptogr.}, 77, 2015.

\bibitem{bernstein2006curve25519}
D.\,J.~Bernstein,
``Curve25519,''
\emph{PKC 2006}, pp.~207--228, 2006.

\bibitem{bernstein2012eddsa}
D.\,J.~Bernstein \emph{et al.},
``High-speed high-security signatures,''
\emph{J.\ Crypt.\ Eng.}, 2(2), 2012.

\bibitem{mcgrew2004aes}
D.~McGrew and J.~Viega,
``GCM,''
\emph{NIST}, 2004.

\bibitem{peikert2016decade}
C.~Peikert,
``A Decade of Lattice Cryptography,''
\emph{FnTTCS}, 10(4), 2016.

\bibitem{regev2005lattices}
O.~Regev,
``On lattices, learning with errors,''
\emph{STOC}, pp.~84--93, 2005.

\bibitem{ajtai1996generating}
M.~Ajtai,
``Generating hard instances of lattice problems,''
\emph{STOC}, pp.~99--108, 1996.

\bibitem{lyubashevsky2010ideal}
V.~Lyubashevsky, C.~Peikert, and O.~Regev,
``On ideal lattices and learning with errors over rings,''
\emph{EUROCRYPT}, 2010.

\bibitem{giacon2018hybrid}
F.~Giacon, F.~Heuer, and B.~Poettering,
``KEM Combiners,''
\emph{PKC 2018}, pp.~190--218, 2018.

\bibitem{bindel2019hybrid}
N.~Bindel \emph{et al.},
``Hybrid Key Encapsulation Mechanisms,''
\emph{PQCrypto 2019}, 2019.

\bibitem{almeida2017jasmin}
J.\,B.~Almeida \emph{et al.},
``Jasmin: High-Assurance and High-Speed Cryptography,''
\emph{CCS}, 2017.

\bibitem{blanchet2008cryptoverif}
B.~Blanchet,
``CryptoVerif: Computationally Sound Mechanized Prover,''
\emph{Dagstuhl Seminar}, 2008.

\bibitem{wood2014ethereum}
G.~Wood,
``Ethereum Yellow Paper,''
2014.

\end{thebibliography}

\appendix

\section{Algorithm OIDs}
\label{app:oids}

\begin{table}[h]
\centering
\begin{tabular}{@{}ll@{}}
\toprule
\textbf{Algorithm} & \textbf{OID} \\
\midrule
ML-KEM-768        & 2.16.840.1.101.3.4.4.2 \\
ML-DSA-65         & 2.16.840.1.101.3.4.3.18 \\
SLH-DSA-SHA2-128s & 2.16.840.1.101.3.4.3.20 \\
\bottomrule
\end{tabular}
\end{table}

\section{Hybrid KEM Test Vectors}
\label{app:vectors}

\begin{lstlisting}[caption={Test Vector Structure}]
type HybridKEMTestVector struct {
    ClassicalSK  [32]byte  // X25519
    ClassicalPK  [32]byte
    PQSK         []byte    // ML-KEM-768
    PQPK         []byte
    Ciphertext   []byte    // 1120 bytes
    SharedSecret [32]byte
}

// Full vectors at:
// github.com/pars-network/pqc-test-vectors
\end{lstlisting}

\section{Migration Checklist}
\label{app:checklist}

\begin{enumerate}
    \item Integrate ML-KEM-768 and ML-DSA-65 libraries
    \item Add hybrid key pair generation to wallet software
    \item Extend prekey bundles with PQ components
    \item Deploy hybrid X3DH in session protocol
    \item Add PQ key registration to L2 account registry
    \item Update bridge contracts for hybrid verification
    \item Enable SLH-DSA for governance key ceremonies
    \item Run interoperability tests across all clients
    \item Publish migration documentation for third-party developers
    \item Conduct security audit of PQ implementations
\end{enumerate}

\end{document}
